\documentclass[10pt]{article}
\usepackage[T1]{fontenc}
\usepackage[utf8]{inputenc}
\usepackage[french]{babel}
\usepackage{amsmath,amssymb,mathrsfs,mathpazo}
\usepackage{geometry}
\geometry{a4paper,margin=2cm}
\usepackage{graphicx,array,tabularx,tabularray,colortbl,multirow,multicol,booktabs,enumitem,pifont,ulem,fancybox,diagbox,cancel}
\usepackage{pstricks,pstricks-add,pst-plot,pst-tree,pst-node,pst-text,pst-func,pst-math,pst-bspline,pst-3dplot,pst-eucl,pst-solides3d}
\let\pscalculate\undefined
\usepackage{pst-bezier}
\usepackage{tikz,pgfplots,tkz-tab}
\usetikzlibrary{arrows,arrows.meta,patterns,calc,positioning,trees,decorations.markings,shapes.geometric,shapes.misc,shapes.arrows,intersections,decorations.pathmorphing,decorations.pathreplacing,backgrounds,shadows,babel}
\pgfplotsset{compat=1.18}
\usepgfplotslibrary{fillbetween}
\usepackage[np]{numprint}
\usepackage{listings,makecell,pas-tableur,twoopt}
\let\barre\undefined
\usepackage[pstricks]{tablvar}
\definecolor{darkgreen}{RGB}{0,100,0}
\providecommand{\red}{\color{red}}
\providecommand{\blue}{\color{blue}}
\providecommand{\green}{\color{green}}
\providecommand{\black}{\color{black}}
\providecommand{\white}{\color{white}}
\providecommand{\gray}{\color{gray}}
\providecommand{\cyan}{\color{cyan}}
\providecommand{\magenta}{\color{magenta}}
\providecommand{\yellow}{\color{yellow}}
\providecommand{\violet}{\color{violet}}
\providecommand{\orange}{\color{orange}}
\providecommand{\euro}{\texteuro}
\providecommand{\EUR}{\texteuro}
\providecommand{\ieme}{\textsuperscript{e}}
\providecommand{\degres}{\ensuremath{^\circ}}
\providecommand{\degree}{\ensuremath{^\circ}}
\providecommand{\vv}[1]{\overrightarrow{#1}}
\providecommand{\Coord}[1]{#1}
\providecommand{\decosix}{}
\providecommand{\starredbullet}{\ensuremath{\bullet}}
\pagestyle{empty}
\setlength{\parindent}{0pt}
\usepackage{adjustbox}\DeclareUnicodeCharacter{2236}{:}\DeclareUnicodeCharacter{2190}{\ensuremath{\leftarrow}}
\providecommand{\R}{\mathbb{R}}
\providecommand{\N}{\mathbb{N}}
\providecommand{\Z}{\mathbb{Z}}
\providecommand{\Q}{\mathbb{Q}}
\providecommand{\D}{\mathbb{D}}
\providecommand{\C}{\mathbb{C}}
\providecommand{\e}{{\mathrm{e}}}
\providecommand{\vect}[1]{\overrightarrow{#1}}
\providecommand{\vectt}[1]{\overrightarrow{\mathrm{#1}}}
\providecommand{\vvt}[1]{\overrightarrow{\mathrm{#1}}}
\providecommand{\np}[2][]{{#2~#1}}
\providecommand{\vv}[1]{{\overrightarrow{#1}}}
\renewcommand{\d}{{\mathrm{d}}}
\providecommand{\degres}{{^\circ}}
\providecommand{\degre}{{^\circ}}
\providecommand{\blue}{\color{blue}}
\providecommand{\red}{\color{red}}
\providecommand{\enskip}{\quad}
\providecommand{\llbracket}{\lbrack\!\lbrack}
\providecommand{\rrbracket}{\rbrack\!\rbrack}
\providecommand{\euro}{{\text{€}}}
\providecommand{\no}{{\text{n°}}}
\providecommand{\interval}[2]{{[#1;#2]}}
\providecommand{\intervFO}[2]{{[#1;#2[}}
\providecommand{\dotfill}{\ldots}

\providecommand{\ds}{\displaystyle}
\providecommand{\textding}[1]{#1}
\providecommand{\Oij}{\(\mathrm{O};\vec{\imath},\vec{\jmath}\)}
\providecommand{\Oijk}{\(\mathrm{O};\vec{\imath},\vec{\jmath},\vec{k}\)}
\providecommand{\Ouv}{\(\mathrm{O};\vec{u},\vec{v}\)}
\providecommand{\intervFF}[2]{\ensuremath{[#1;#2]}}
\renewcommand{\euro}{\ifmmode\text{\texteuro}\else\texteuro\fi}
\DeclareUnicodeCharacter{00B0}{\ifmmode^{\circ}\else\textdegree\fi}
\DeclareUnicodeCharacter{20AC}{\ifmmode\text{\texteuro}\else\texteuro\fi}
% Énoncés : Éducation nationale. Transcription et corrigés : APMEP, Denis Vergès et contributeurs. Découpage et classement : CoopMaths, Sébastien Lozano, Éric Elter et contributeurs. Adaptation de la présentation : Maths Explorer.
% CC BY-SA : https://coopmaths.fr/www/about/
% https://coopmaths.fr/alea/static/bac/2025/tex/bac_2025_06_sujet1_etrangers_4.tex
\begin{document}
\noindent\textit{APMEP / CoopMaths. CC BY-SA. Présentation adaptée par Maths Explorer.}\par\medskip
\subsection*{Partie A}

On considère l'équation différentielle 

\[(E_1) :\quad  y' + 0{,}48y = \dfrac{1}{250},\]

où $y$ est une fonction de la variable $t$ appartenant à l'intervalle $[0~;~ +\infty[$.

\begin{enumerate}
\item On considère la fonction constante $h$ définie sur l'intervalle $[0~;~ +\infty[$ par $h(t) = \dfrac{1}{120}$.

Montrer que la fonction $h$ est solution de l'équation différentielle $(E_1)$.
\item Donner la forme générale des solutions de l'équation différentielle $y' + 0,48y = 0$.
\item En déduire l'ensemble des solutions de l'équation différentielle $(E_1)$.
\end{enumerate}

\subsection*{Partie B}

On s'intéresse à présent à l'évolution d'une population de bactéries dans un milieu de culture.

À un instant $t = 0$, on introduit une population initiale de \np{30000} bactéries dans le milieu. On note $p(t)$ la quantité de bactéries, exprimée en millier d'individus, présente dans le milieu après un temps $t$, exprimé en heure.

On a donc $p(0) = 30$.

On admet que la fonction $p$ définie sur l'intervalle $[0~;~ +\infty[$ est dérivable, strictement positive sur cet intervalle et qu'elle est solution de l'équation différentielle $(E_2)$ :

\[p' = \dfrac{1}{250}p \times (120 - p)\]

Soit $y$ la fonction strictement positive sur l'intervalle $[0~;~ +\infty[$ telle que, pour tout $t$ appartenant à l'intervalle $[0~;~ +\infty[$ , on a $p(t) = \dfrac{1}{y(t)}$.

\medskip

\begin{enumerate}
\item Montrer que si $p$ est solution de l'équation différentielle $(E_2)$, alors $y$ est solution de l'équation différentielle $(E_1)$ : \quad $y' + 0,48y = \dfrac{1}{250}$.
\item On admet réciproquement que, si $y$ est une solution strictement positive de l'équation différentielle $(E_1)$, alors $p = \dfrac 1y$ est solution de l'équation différentielle $(E_2)$.

Montrer que, pour tout $t$ appartenant à l'intervalle $[0~;~ +\infty[$ , on a :

\[p(t) = \dfrac{120}{1 + K\e^{-0,48t}}~ \text{avec}\: K\: \text{une constante réelle}.\] 
\item En utilisant la condition initiale, déterminer la valeur de $K$.
\item Déterminer $\displaystyle\lim_{t \to + \infty} p(t)$. En donner une interprétation dans le contexte de l'exercice.
\item Déterminer le temps nécessaire pour que la population de bactéries dépasse \np{60000}~individus.

\emph{On donnera le résultat sous la forme d'une valeur arrondie exprimée en heures et minutes.}
\end{enumerate}

\end{document}
